\documentclass{article}
\input{~/studies/preamble}

\title{Probabilistic Principal Component Analysis}
\author{Hyoung Joon, Kim}
\date{\today}

\begin{document}
\maketitle
\tableofcontents
\newpage

\section{Introduction}
    \textbf{Probabilistic principal component analysis} (pPCA) is the probabilistic version of PCA
    using continuous latent variables. The advantage of pPCA over PCA is:
    \begin{itemize}
        \item Use EM algorithm for PCA that is computationally efficient when only a few eigenvectors are required;
        \item Direct comparison with other probabilistic density models, using likelihood function;
        \item Can be applied to classification problems;
        \item Can be run generatively to provide samples from the distribution.
    \end{itemize}
\section{Main Ideas}
    \subsection{Construction of pPCA}
        pPCA is a simple example of linear-Gaussian framework where all marginal and conditional distributions are Gaussian.
        We define a Gaussian prior distribution $p(\textbf{z})$ over latent variable $\textbf{z}$ as $p(\textbf{z})=\mathcal{N}(\textbf{z}|\textbf{0},\textbf{I})$.
        Then, we define the conditional distribution of the observed variable \textbf{x} conditioned over \textbf{z} as
        $p(\textbf{x}|\textbf{z})=\mathcal{N}(\textbf{x}|\textbf{Wz}+\bm\mu,\sigma^2\textbf{I})$.
        This is equal to $\tb x = \tb{Wz} + +\bm\mu + \bm\epsilon$ where $\bm\epsilon$ is a Gaussian noise with mean 0 and isotropic covariance.

        \begin{defn}{Probabilistic Principal Component Analysis}{ppca}
            Let \textbf{z} be a latent variable of $M$ dimension, and let \textbf{x} be observed data. Then, the PPCA is a probabilistic model with following
            probability distributions;
            \begin{align*}
                p(\textbf{z})&=\mathcal{N}(\textbf{0},\textbf{I}) \\
                p(\textbf{x}|\textbf{z},\bm\theta)&=\mathcal{N}(\textbf{x}|\textbf{Wz}+\bm\mu,\sigma^2\textbf{I})
            \end{align*}
            where \tb W is a $D \times M$ weight matrix.
            Then the marginal distribution of \textbf{x} is given as
            \[
                p(\textbf{x}|\bm\theta)=\mathcal{N}(\textbf{x}|\bm\mu,\textbf{C})
            \]
            where \textbf{C} is defined by $\textbf{C}=\textbf{W}\textbf{W}^\top + \sigma^2\textbf{I}$.

            The posterior distribution of latent variable \textbf{z} is given as
            \[
                p(\textbf{z}|\textbf{x},\bm\theta)=\mathcal{N}(\textbf{z}|\tb{M}^{-1}\tb{W}^\top(\tb{x}-\bm\mu),\sigma^2\tb{M}^{-1})
            \]
            where \tb{M} is defined by $\tb{M}=\tb{W}^\top\tb{W}+\sigma^2\tb{I}$
        \end{defn}

        We can notice that $\sigma^2 \rightarrow 0$ leads to the original PCA.

    \subsection{Maximum Likelihood for pPCA}
        \begin{thm}{Maximum Likelihood for pPCA}{maxliPPCA}
            The log-likelihood of pPCA is given as
            \begin{align*}
                \log{p(\textbf{X}|\bm\mu\tb{W},\sigma^2)} &= \sum_{n=1}^{N}\log{p(\tb{x}_n|\tb W, \bm\mu,\sigma^2)} \\
                &= -\frac{ND}{2}\log{2\pi}-\frac{N}{2}\log{|\tb C|}-\frac{1}{2}\sum_{n=1}^{N}(\tb x_n -\bm\mu)^\top\tb{C}^{-1}(\tb x_n - \bm\mu)
            \end{align*}
            The MLE of each parameters are given as:
            \begin{align*}
                \bm\mu &= \overline{\tb x} \\
                \tb W &= \tb U_M(\tb L_M -\sigma^2\tb I)^{1/2}\tb R \\
                \sigma^2 &=\frac{1}{D-M}\sum_{i=M+1}^{D}\lambda_i
            \end{align*}
            where $\tb U_M$ is a $D \times M$ matrix whose columns are given by $M$ largest eigenvectors of the data covariance matrix $\tb S$.
            The $\tb L_M$ has elements given by the corresponding eigenvalues $\lambda_i$, and \tb R is arbitrary $M \times M$ ortogonal matrix.
            The summation of MLE of $\sigma^2$ is about the leftover eigenvalues of eigenvectors not contained in $\tb U_M$. 
        \end{thm}
        Notice that the likelihood function is invariant over the rotation of weight matrix $\tb W$,
        since $\tb C = \tilde{\textbf{W}}\tilde{\textbf{W}}^\top + \sigma^2\textbf{I}$ where $\tilde{\tb W}=\tb{WR}$ with arbitrary orthogonal matrix $\tb R$.

    \subsection{EM for pPCA}
        \subsubsection{Expecation Maximization}
            \begin{algorithm}[H]
                \DontPrintSemicolon
                \SetKwInOut{Input}{Input}
                \SetKwInOut{Output}{Output}
                \SetArgSty{textnormal}

                \caption{EM algorithm for pPCA}
                
                \Input{dataset $\tb X$ with dimensionality $N \times D$\newline
                initial parameter values $\tb W$, $\sigma^2$}
                \Output{final parameter values $\tb W$, $\sigma^2$}

                \BlankLine
                \hrule
                \tcp{initialize}
                \Repeat{convergence}
                {
                    \tcp{E step}
                    $\tb M \leftarrow \tb{W}^\top\tb{W}+\sigma^2\tb{I}$\;
                    $\mathbb{E}[\textbf{z}_n] \leftarrow \tb M^{-1}\tb W^\top(\tb x_n - \bar{\tb{x}})$\;
                    $\mathbb{E}[\tb z_n \tb z^\top] \leftarrow \sigma^2\tb M^{-1} + \mathbb{E}[\tb z_n]\mathbb{E}[\tb z_n]^\top$\;
                    \tcp{M step}
                    $\tb W_{\text{new}} \leftarrow \left[ \sum_{n=1}^{N}(\tb x_n - \overline{\tb x})\mathbb{E}[\tb z_n]^\top \right] \left[ \sum_{n=1}^{N}\mathbb{E}[\tb z_n\tb z_n^\top] \right]^{-1}$\;
                    $\sigma_{\text{new}}^2 \leftarrow \frac{1}{ND}\sum_{n=1}^{N}\left\{ \|\tb x_n \overline{\tb x}\|^2 -2\mathbb{E}[\tb z_n]^\top\tb W_{\text{new}}^\top(\tb x_n -\overline{\tb x})
                    + \tb{Tr}(\mathbb{E}[\tb z_n \tb z_n^\top]\tb W_new^\top \tb W_new)\right\}$\;
                    $\tb W, \sigma^2 \leftarrow \tb W_{\text{new}}, \sigma^2_{\text{new}}$\;
                }
                \Return{$\tb W$, $\sigma^2$}
                \end{algorithm}
        \subsubsection{ELBO}
            EM for continuous latent variables are derived in the same way as EM for discrete latent variable. See \href{run:../inference/em_algorithm/em_algorithm.pdf}{this document}

\section{Extensions of pPCA}
    \subsection{Factor Analysis}
        In probabilistic PCA, the covariance matrix of latent variable $\tb z$ is given as an identity matrix. If we changed the covariance matrix from identity matrix to diagonal matrix,
        we obtain \textbf{Factor analysis} model. We assume that the data are linear combination of latent variables plus some Gaussian noise, i.e., $\tb x = \tb{Wz} + \bm\mu + \bm\epsilon$.
        The main difference from pPCA is, that the noise's variance differes between features, i.e., $\bm\epsilon \sim \mathcal{N}(\tb 0, \bm\Psi)$ where $\bm\Psi$ is diagonal matrix.
        The conditional probability of data is given as
        \[
        p(\tb x| \tb z)=\mathcal{N}(\tb x|\tb W \tb z + \bm\mu, \bm\Psi)
        \]
        where $\bm\Psi$ is a diagonal noise matrix. The marginal distribution is then given as
        \[
        p(\tb x)=\mathcal{N}(\tb x|\bm\mu,\tb C)
        \]
        where $\tb C=\tb W\tb W^\top + \Psi$. This model finds out the correlation between variables, assuming that the correlation stems from common latent variable.
        Since the noises can differ between features, it can capture the correlation between features well, independent of noises.
    \subsection{Independent Component Analysis}
        One of the generalizations of probabilistc PCA assumes that the observed variables are related linearly to the latent variables. The difference is that
        the latent variables are not Gaussian anymore. \textbf{Independent component analysis} is an important class of such models, which assumes that the latent variables
        factorizes completely, i.e., $p(\tb z) = \prod_{j=1}^{M}p(z_j)$.

        ICA is used to analyze the idependent components of given data, as we can see from its name. The situation which explains this well is the 'Cocktail Party Problem'.
        ICA assumes that the original data are the linear combinations of the sources, i.e., $\tb x = \tb A\tb s$, and assumes that each element of $\tb s$ is independent, i.e., $p(\tb s)=\prod_{i=1}^{M}p_i(s_i)$.
        The objective of ICA is to find out the best decomposition matrix $\tb W$ such that $\tb s = \tb W \tb x$. To do so, we maximize the likelihood of dataset, about $\tb W$, which is given as
        \[
            p(\tb x) = \left( \prod_{i=1}^{n}p_i(\tb w_i^\top \tb x) \right)\|\tb W\|
        \]
        ICA requires the distribution of latent variables to be non-Gaussian, since it won't be able to distinguish between two different choices for latent variables
        if these differ simply by a rotation in latent space, as we saw in Theorem \ref{thm:maxliPPCA}.
        Since the likelihood function produces the same probability even when the latent variable is rotated, there are infinite choices for maximum likelihood. This means that
        the original sources are impossible to extract since $\tb s = \tb W \tb x$ will produce different $\tb s$ for each rotated $\tb W$. To solve this, ICA
        uses non-Gaussian distribution, mainly hyperbolic secant, $p_i(s_i)=\frac{1}{\pi\text{cosh}(s_i)}$.

\end{document}